\newpage
\section{Lecture-09}
\begin{proposition}
  If the function $f:U\subseteq \R^n \to \R$ is differentiable at $\vec a$, then the partial derivatives 
  $\frac{\partial f}{\partial x_i}(\vec a)$ exists for all $i=1,2,\ldots,n$ and $Df(\vec a)$ is given by the row vector
  \[
  Df(\vec a) = \begin{pmatrix}
  \frac{\partial f}{\partial x_1}(\vec a) & \frac{\partial f}{\partial x_2}(\vec a) & \cdots & \frac{\partial f}{\partial x_n}(\vec a)
  \end{pmatrix}.
  \] 
\end{proposition}
\begin{proof}
  $f$ is differentiable at $\vec a$ means that there exists a linear transformation $L\equiv Df(\vec a):\R^n \to \R$ such that
  \[
  f(\vec a+\vec h)=f(\vec a)+L(\vec h)+o(\|\vec h\|).
  \]
  where $\lim_{\|h\|\rightarrow 0} \frac{o(\|\vec h\|)}{\|\vec h\|}=0$. Suppose, 
  $$L(\vec h)=L\begin{pmatrix}h_1\\h_2\\\vdots\\h_n\end{pmatrix}=\begin{pmatrix}
  c_1 & c_2 & \cdots & c_n
  \end{pmatrix}\begin{pmatrix}h_1\\h_2\\\vdots\\h_n\end{pmatrix}=c_1h_1+c_2h_2+\cdots+c_nh_n.$$
  Now, we want to impose the definition in each of the coordinate directions. For $i=1,2,\ldots,n$, let $\vec h = t \vec e_i$ where $\vec e_i$ is the $i$-th standard basis vector. Then
  \begin{align*}
  f(\vec a+t\vec e_i)&=f(\vec a)+L(t\vec e_i)+o(\|t\vec e_i\|)\\
  f(\vec a+t\vec e_i) - f(\vec a) &= t L(\vec e_i) + o(|t|)\\
  \frac{f(\vec a+t\vec e_i) - f(\vec a)}{t} &= L(\vec e_i) + \frac{o(|t|)}{t}\\
  \lim_{t \to 0} \frac{f(\vec a+t\vec e_i) - f(\vec a)}{t} &= L(\vec e_i) + \lim_{t \to 0} \frac{o(|t|)}{t}\\
  \frac{\partial f}{\partial x_i}(\vec a) &= L(\vec e_i) = c_i.
  \end{align*}
\end{proof}
\begin{proposition}
  \label{total_derivative_matrix}
  If the function $f:U\subseteq \R^n \to \R^m$ is differentiable at $\vec a$, then the partial derivatives 
  $\frac{\partial f^j}{\partial x_i}(\vec a)$ exists for all $i=1,2,\ldots,n$, $j=1,2,\ldots,m$ and $Df(\vec a)$ is given by the matrix
  \[
  Df(\vec a) = \begin{pmatrix}
  \frac{\partial f^1}{\partial x_1}(\vec a) & \frac{\partial f^1}{\partial x_2}(\vec a) & \cdots & \frac{\partial f^1}{\partial x_n}(\vec a)\\
  \frac{\partial f^2}{\partial x_1}(\vec a) & \frac{\partial f^2}{\partial x_2}(\vec a) & \cdots & \frac{\partial f^2}{\partial x_n}(\vec a)\\
  \vdots & \vdots & \ddots & \vdots \\
  \frac{\partial f^m}{\partial x_1}(\vec a) & \frac{\partial f^m}{\partial x_2}(\vec a) & \cdots & \frac{\partial f^m}{\partial x_n}(\vec a)
  \end{pmatrix}.
  \]
\end{proposition} 
% =========================================
\begin{proposition}{(Optional)}
Let $f:\R^{k}\to\R^{n}$ be a differentiable map. For each $i\in\{1,\dots,n\}$ let
\[
\pi^i:\R^n\to\R,\qquad \pi^i(a^1,\dots,a^n)=a^i
\]
be the $i$-th coordinate projection, and define the $i$-th component function of $f$ by
\[
f^i := \pi^i\circ f:\R^k\to\R.
\]
Then for every $a\in\R^k$,
\[
D(f^i)(a) = \pi^i\circ Df(a).
\]
Equivalently, as linear maps $\R^k\to\R$,
\[
D(f^i)(a)[h] = \pi^i\big(Df(a)[h]\big)\qquad\text{for all }h\in\R^k.
\]
\end{proposition}
\begin{proof}
  See the proof from \url{https://math.stackexchange.com/a/4214572/1167703}.
\end{proof}
% \begin{proof}
% By definition $f^i=\pi^i\circ f$. Since $f$ is differentiable at $a$ and $\pi^i$ is differentiable everywhere, the chain rule gives
% \[
% D(f^i)(a)=D(\pi^i\circ f)(a)=D\pi^i(f(a))\circ Df(a).
% \]
% Now observe that $\pi^i:\R^n\to\R$ is a linear transformation. A basic fact about differentiability is that if $L:\R^n\to\R^m$ is linear, then
% \[
% DL(x)=L \quad\text{for every }x\in\R^n.
% \]
% Applying this to $L=\pi^i$ yields $D\pi^i(f(a))=\pi^i$. Therefore,
% \[
% D(f^i)(a)=\pi^i\circ Df(a),
% \]
% as claimed. Writing this identity on an increment $h\in\R^k$ gives
% \[
% D(f^i)(a)[h]=(\pi^i\circ Df(a))[h]=\pi^i(Df(a)[h]).
% \]
% \end{proof}

% \begin{remark}[What the notation means]
% It is \emph{not} correct to write $\pi^i(f)=f^i$ without explanation: $\pi^i$ is a function $\R^n\to\R$ whose input is a point
% $(a^1,\dots,a^n)$, not a function. The correct definition is
% \[
% f^i := \pi^i\circ f.
% \]
% \end{remark}

% %------------------------------------------------------------
% % Example 1: Spivak-style D as linear map (no matrices needed)
% %------------------------------------------------------------
% \begin{example}[A concrete computation with $D$ as a linear map]
% Let $f:\R\to\R^2$ be defined by
% \[
% f(x)=(x^2,x).
% \]
% Then for each $a\in\R$, the total derivative $Df(a):\R\to\R^2$ is the linear map determined by
% \[
% Df(a)[h] = (2ah,\,h)\qquad (h\in\R).
% \]
% Let $\pi^1:\R^2\to\R$ be the projection $\pi^1(u,v)=u$. Then $f^1=\pi^1\circ f$ satisfies
% \[
% f^1(x)=\pi^1(x^2,x)=x^2.
% \]
% By the proposition,
% \[
% Df^1(a)[h] = \pi^1(Df(a)[h])=\pi^1(2ah,h)=2ah,
% \]
% so indeed $(f^1)'(a)=2a$.
% \end{example}

% %------------------------------------------------------------
% % Example 2: Matrix viewpoint (prime notation)
% %------------------------------------------------------------
% \begin{example}[The same example using matrices]
% Again let $f(x)=(x^2,x)$. Using standard bases, the Jacobian (matrix representation of $Df(a)$) is
% \[
% f'(a)=
% \begin{pmatrix}
% 2a\\[0.3ex]
% 1
% \end{pmatrix}
% \qquad \text{(a $2\times 1$ matrix).}
% \]
% The projection $\pi^1:\R^2\to\R$ is linear, and its matrix (standard basis in $\R^2$ and $\{1\}$ in $\R$) is
% \[
% [\pi^1]_{\text{std}}=
% \begin{pmatrix}
% 1 & 0
% \end{pmatrix}
% \qquad \text{(a $1\times 2$ matrix).}
% \]
% The chain rule in matrix form gives
% \[
% (f^1)'(a) = (\pi^1)'(f(a))\, f'(a) = [\pi^1]_{\text{std}}\, f'(a)
% =
% \begin{pmatrix}
% 1 & 0
% \end{pmatrix}
% \begin{pmatrix}
% 2a\\[0.3ex]
% 1
% \end{pmatrix}
% =2a.
% \]
% \end{example}

% %------------------------------------------------------------
% % Example 3: A higher-dimensional example to show the pattern
% %------------------------------------------------------------
% \begin{example}[General pattern in $\R^k\to\R^n$]
% Let $f:\R^2\to\R^3$ be
% \[
% f(x,y)=(x^2+ y,\; xy,\; e^x).
% \]
% Then $f^2=\pi^2\circ f$ is the second component:
% \[
% f^2(x,y)=xy.
% \]
% The total derivative $Df(a,b):\R^2\to\R^3$ is given (on $(h_1,h_2)\in\R^2$) by
% \[
% Df(a,b)[h_1,h_2]=\big(2a\,h_1+h_2,\; b\,h_1+a\,h_2,\; e^a h_1\big).
% \]
% Applying the proposition with $\pi^2(u_1,u_2,u_3)=u_2$,
% \[
% Df^2(a,b)[h_1,h_2]=\pi^2(Df(a,b)[h_1,h_2])=b\,h_1+a\,h_2,
% \]
% which matches the usual gradient computation for $xy$.
% \end{example}
% =========================================
\begin{proposition}
    If $f: A \subseteq \R^n \rightarrow \R^m$ is differentiable at $x_0$, then $f$ is continuous at $x_0$.
\end{proposition}
\begin{proof}
    By definition of the derivative of $f$ we know that the limit
$$
\lim_{h \rightarrow 0} \frac{|| f(x_0+h) - f(x_0) + D h || }{||h||} = 0,
$$
for some linear transformation $D$. Therefore since this limit exists, the following is true:
$$
0 = \left( \lim_{h \rightarrow 0} || h ||\right)\left( \lim_{h \rightarrow 0} \frac{|| f(x_0+h) - f(x_0) + D h || }{||h||} \right) = \lim_{h \rightarrow 0} ||f(x_0+h) - f(x_0) + Dh ||.
$$
By the triangle inequality,
$$
\lim_{h \rightarrow 0}||f(x_0+h) - f(x_0) || \leq \lim_{h \rightarrow 0}||f(x_0+h) - f(x_0) +Dh || + \lim_{h \rightarrow 0}||Dh || =\lim_{h \rightarrow 0}||Dh ||.
$$
But since $D$ is linear, it is continuous and $D0=0$, so the lattermost limit is just $0$. Thus $0$ is an upper bound for the limit, but since the norm is always non-negative, this implies that the limit is in fact $0$.
\end{proof}
% =======================================
\begin{example}
\[
f(x,y,z)=x^2+y^2+z^2.
\]
Take 
\[
\vec a=(1,1,1).
\]
\[
Df(x,y,z)=(2x,2y,2z).
\]
\[
Df(1,1,1)=(2,2,2).
\]
Thus, for $\vec v=(v_1,v_2,v_3)$,
\[
Df(1,1,1)(\vec v)
=
2v_1+2v_2+2v_3.
\]
\[
D_{\vec v}f(1,1,1)
=
\lim_{t\to0}
\frac{(1+tv_1)^2+(1+tv_2)^2+(1+tv_3)^2-3}{t}.
\]
Expanding,
\[
(1+tv_1)^2+(1+tv_2)^2+(1+tv_3)^2
=
3+2t(v_1+v_2+v_3)+O(t^2).
\]
Therefore,
\[
D_{\vec v}f(1,1,1)
=
\lim_{t\to0}
\frac{3+2t(v_1+v_2+v_3)+O(t^2)-3}{t}
=
2(v_1+v_2+v_3).
\]
Hence,
\[
\boxed{D_{\vec v}f(\vec a)=Df(\vec a)(\vec v)}
\]
\end{example}
~\\
We can see why it works in general.
If $f$ is differentiable at $\vec a$, then
\begin{align*}
f(\vec a+\vec h)&=f(\vec a)+Df(\vec a)(\vec h)+o(|\vec h|)\\
\stackrel{\vec h=t\vec v}{\implies}\frac{f(\vec a+t\vec v)-f(\vec a)}{t}&=Df(\vec a)(\vec v)+\frac{o(t|\vec v|)}{t}.
\end{align*}
As $t\to0$, the error term goes to $0$, giving
\[
D_{\vec v}f(\vec a)=Df(\vec a)(\vec v).
\]
For more rigorous proof, see the following proposition.
\begin{proposition}
  When $f$ is differentiable at $\vec a$, for any $\vec v$, 
  the directional derivative of $f$ at $\vec a$ in the direction $\vec v$ exists and is given by
  \[
  D_{\vec v}f(\vec a) = Df(\vec a)(\vec v).
  \]    
\end{proposition}
\begin{proof}
  See the proof from Theodore Shifrin \cite{shifrin} (Chapter 3, section 2, proposition 2.3). 
\end{proof}
\begin{example}
  Let $f:\mathbb R^2\to \mathbb R^2$ be defined by,
  \[
  f\begin{pmatrix}
    x\\
    y
  \end{pmatrix}=\begin{pmatrix}
    xy\\
    x^2-y^2
  \end{pmatrix}
  \]
  At the point $\vec a=(1,1)$, at what rate $f$ varying in the
  direction $\vec v=(1,0)$?
\end{example}
The derivative of $f$ at $\vec a$ is,
\[
Df(\vec x)=\begin{pmatrix}
y & x\\
2x & -2y
\end{pmatrix}.
\]
Thus, 
\[
D_{\vec v}f(\vec a)=Df(\vec a)(\vec v)=\begin{pmatrix}
1 & 1\\
2 & -2
\end{pmatrix}\begin{pmatrix}
2\\
1
\end{pmatrix}=\begin{pmatrix}
3\\
2
\end{pmatrix}.
\]
See, how the directional derivative is just the application of the total derivative to the direction vector.
%------------------------------------------------------------
% Properties of directional derivatives (LaTeX snippet)
%------------------------------------------------------------
% ~\\
% Let $f,g:\mathbb{R}^n\to\mathbb{R}$, let $p\in\mathbb{R}^n$, and let $v\in\mathbb{R}^n$.
% The \emph{directional derivative} of $f$ at $p$ in the direction $v$ is
% \[
% D_v f(p)\;:=\;\left.\frac{d}{dt}\right|_{t=0}\, f(p+t\,v),
% \]
% whenever the limit exists. Because,
% \begin{align*}
%   \frac{d}{dt}f(p+tv) &= \lim_{h\to0}\frac{f(p+(t+h)v)-f(p+tv)}{h}\\
%   \left.\frac{d}{dt}f(p+tv)\right|_{t=0} &= \lim_{h\to0}\frac{f(p+hv)-f(p)}{h}\\
%   &= D_v f(p).
% \end{align*} 
% \bigskip
% \noindent\textbf{(1) Linearity in the function.}
% If $D_v f(p)$ and $D_v g(p)$ exist, then for all scalars $\alpha,\beta\in\mathbb{R}$,
% \[
% D_v(\alpha f+\beta g)(p)=\alpha\,D_v f(p)+\beta\,D_v g(p).
% \]
% \medskip
% \noindent\textbf{(2) Product rule (Leibniz rule).}
% If $D_v f(p)$ and $D_v g(p)$ exist, then
% \[
% D_v(fg)(p)=f(p)\,D_v g(p)+g(p)\,D_v f(p).
% \]
% \medskip
% \noindent\textbf{(3) Quotient rule.}
% If $g(p)\neq 0$ and $D_v f(p),D_v g(p)$ exist, then
% \[
% D_v\!\left(\frac{f}{g}\right)(p)=\frac{g(p)\,D_v f(p)-f(p)\,D_v g(p)}{(g(p))^2}.
% \]
% \medskip
% \noindent\textbf{(4) Chain rule (one-variable outer function).}
% If $\phi:\mathbb{R}\to\mathbb{R}$ is differentiable at $f(p)$ and $D_v f(p)$ exists, then
% \[
% D_v(\phi\circ f)(p)=\phi'(f(p))\,D_v f(p).
% \]
% \medskip
% \noindent\textbf{(5) Scaling in the direction.}
% For any scalar $c\in\mathbb{R}$ (whenever $D_v f(p)$ exists),
% \[
% D_{c v} f(p)=c\,D_v f(p).
% \]
% \medskip
% \noindent\textbf{(6) Additivity/linearity in the direction (requires differentiability).}
% If $f$ is differentiable at $p$, then for all $v,w\in\mathbb{R}^n$ and all $a,b\in\mathbb{R}$,
% \[
% D_{v+w} f(p)=D_v f(p)+D_w f(p),
% \qquad\text{and more generally}\qquad
% D_{a v+b w} f(p)=a\,D_v f(p)+b\,D_w f(p).
% \]


\begin{proposition}
  If $f: U \subseteq \R^n \rightarrow \R^m$ has continuous partial derivatives, then $f$ is differentiable.  
\end{proposition}
\begin{proof}
  See the proof from Theodore Shifrin \cite{shifrin} (Chapter 3, section 2, proposition 2.4). 
\end{proof}
% continuous partial DNI differentiability: https://mathinsight.org/differentiable_function_discontinuous_partial_derivatives
\subsection{Jacobian Matrix}
\begin{definition}
For a function $f:\mathbb{R}^n \to \mathbb{R}^m$, written as
\[
f(\vec x)=f\begin{pmatrix}
x_1\\
x_2\\
\vdots\\
x_n
\end{pmatrix}=\begin{pmatrix}
f_1(x_1,\dots,x_n)\\
f_2(x_1,\dots,x_n)\\
\vdots\\
f_m(x_1,\dots,x_n)
\end{pmatrix}=\begin{pmatrix}
f_1(\vec x)\\
f_2(\vec x)\\
\vdots\\
f_m(\vec x)
\end{pmatrix},
\]
the Jacobian matrix of $f$ at a point $x \in \mathbb{R}^n$ is the $(m \times n)$ matrix of first-order partial derivatives:

\[
J_f(x) = \begin{bmatrix}
\frac{\partial f_1}{\partial x_1}(x) & \frac{\partial f_1}{\partial x_2}(x) & \cdots & \frac{\partial f_1}{\partial x_n}(x) \\
\frac{\partial f_2}{\partial x_1}(x) & \frac{\partial f_2}{\partial x_2}(x) & \cdots & \frac{\partial f_2}{\partial x_n}(x) \\
\vdots & \vdots & \ddots & \vdots \\
\frac{\partial f_m}{\partial x_1}(x) & \frac{\partial f_m}{\partial x_2}(x) & \cdots & \frac{\partial f_m}{\partial x_n}(x)
\end{bmatrix}.
\]
Its $(i,j)$-entry is
\[
(J_f(x))_{ij}=\frac{\partial f_i}{\partial x_j}(x).
\]
It represents the best linear approximation of $f$ near $x$. If $m=1$, the Jacobian is just the row vector of partial derivatives, closely related to the gradient.
\end{definition}
\begin{theorem}
  If $f$ is differentiable at $\vec a$, then all the partial derivatives of $f$ at $\vec a$ exist and the matrix representing $Df(\vec a)$ is the Jacobian matrix $J_f(\vec a)$.
\end{theorem}
\begin{proof}
  We already proved this in Proposition \ref{total_derivative_matrix}.
\end{proof}
\begin{remark}
  The gradient of a scalar-valued function $f:\R^n\to\R$ is the row vector of partial derivatives, which is the Jacobian matrix in this case.
  \[
    \operatorname{grad}f(\vec a)=\nabla f(\vec a)=\begin{pmatrix}
    \frac{\partial f}{\partial x_1}(\vec a)\\
    \frac{\partial f}{\partial x_2}(\vec a)\\
    \vdots\\
    \frac{\partial f}{\partial x_n}(\vec a)
    \end{pmatrix}.
  \]
  In other words, $\operatorname{grad}f(\vec a)=[Df(\vec a)]^T=\begin{pmatrix}
    \frac{\partial f}{\partial x_1}(\vec a) & \frac{\partial f}{\partial x_2}(\vec a) & \cdots & \frac{\partial f}{\partial x_n}(\vec a)
  \end{pmatrix}^T$. The transpose is not as innocent as it looks. geometrically, it is only defined if we are in a space with an inner product. When there
  is no nature inner product, the vector $\vec \nabla f(\vec a)$ is not natural, unlike the derivative which is a linear transformation depends on no such choices.    
\end{remark}
